\section{隐函数与参数方程微分法}\label{sec:4.3}

\subsection{隐函数微分法}

前面讲的函数关系，都是通过 $y=f(x)$ 这样的形式给出的，其中 $f(x)$ 是 $x$ 的某种算式，我们称用这种形式表示的函数为显函数．如
\[
y=\sin x+\cos 2x,\qquad y=e^{x}+x^{2}.
\]
有些函数，自变量与因变量之间的对应法则是由一个方程确定的．例如
\[
x^{2}+y^{2}=1,\qquad\frac{x^{2}}{a^{2}}+\frac{y^{2}}{b^{2}}=1.
\]
一般地，对于方程
\[
F(x,y)=0,
\]
若存在集合 $X$，对任意 $x\in X$，存在唯一确定的 $y\in\mathbb{R}$，使得点对 $(x,y)$ 满足上述方程，则称方程 $F(x,y)=0$ 确定了隐函数，记为
\[
y=f(x),\qquad x\in X.
\]
这时有
\[
F(x,f(x))\equiv0,\qquad x\in X.
\]

\begin{example}\label{ex:4-3-1}
方程 $x^{2}+y^{2}=1$ 可以确定隐函数
\[
y=\sqrt{1-x^{2}},\qquad x\in[-1,1];
\]
和
\[
y=-\sqrt{1-x^{2}},\qquad x\in[-1,1];
\]
它们都是连续函数．但也可以确定隐函数
\[
y=\begin{cases}\sqrt{1-x^{2}}, & x\in[0,1],\\ -\sqrt{1-x^{2}}, & x\in[-1,0),\end{cases}
\]
它在 $x=0$ 点不连续．
\end{example}

这个例子中的函数 $y=f(x)$ 是可以从方程``解''出来的．也有些隐函数，例如由下面方程确定的隐函数
\[
y-x-\frac{1}{2}\sin y=0,
\]
它不能简单地解出来．但以后我们会知道，对每个 $x$，的确有且只有一个 $y$ 与之对应，因而它确定 $y$ 作为 $x$ 的函数，不过不能写成 $y$ 作为 $x$ 的初等函数．当然，并非任一方程都可确定隐函数，例如，方程
\[
x^{2}+y^{2}+5=0
\]
在实数系里就不可能确定任何隐函数．

以下我们假定在一定条件下，方程
\[
F(x,y)=0
\]
可以确定隐函数 $y=f(x)$，并且是可导的．在这个前提下，我们通过举例，给出隐函数的微商法．重要的是，在计算过程中，我们无需把 $y$ 解为 $x$ 的显函数．

\begin{example}\label{ex:4-3-2}
由方程 $x^{2}+y^{2}=1$ 确定隐函数 $y=f(x)$，求 $\dfrac{\mathrm{d}y}{\mathrm{d}x}$．

解 将 $y=f(x)$ 代入方程，则方程为恒等式，即有
\[
x^{2}+\left(f(x)\right)^{2}\equiv1.
\]
将 $(f(x))^{2}$ 视为复合函数，在恒等式两边对 $x$ 求导，得恒等式
\[
2x+2f(x)f'(x)\equiv0.
\]
解得
\[
y'=f'(x)=-\frac{x}{f(x)}=-\frac{x}{y}.
\]
也可以利用微分运算求隐函数的导数．在方程
\[
x^{2}+y^{2}=1
\]
两边求微分得
\[
2x\mathrm{d}x+2y\mathrm{d}y=0\cdot\mathrm{d}x.
\]
解得
\[
y'=\frac{\mathrm{d}y}{\mathrm{d}x}=-\frac{x}{y}.
\]
\end{example}

\begin{example}\label{ex:4-3-3}
已知 $y-x-\dfrac{1}{2}\sin y=0$，求 $\dfrac{\mathrm{d}y}{\mathrm{d}x}$．

解 在方程两边对 $x$ 求导，并注意 $y$ 是 $x$ 的函数，得
\[
y'-1-\frac{1}{2}\cos y\cdot y'=0,
\]
解得
\[
y'=\frac{1}{1-\dfrac{1}{2}\cos y}.
\]
\end{example}

\begin{example}\label{ex:4-3-4}
已知 $x^{y}=y^{x}$ ($x>0,y>0$)，求 $\dfrac{\mathrm{d}y}{\mathrm{d}x}$．

解 在方程两边取对数得
\[
y\ln x=x\ln y,
\]
两边对 $x$ 求导，并注意 $y$ 是 $x$ 的函数，得
\[
y'\ln x+\frac{1}{x}y=\ln y+\frac{x}{y}y',
\]
解得
\[
y'=\frac{xy\ln y-y^{2}}{xy\ln x-x^{2}}.
\]
\end{example}

\subsection{参数方程微分法}

在解析几何中，常用参数方程表示曲线．例如椭圆的参数方程为
\[
\begin{cases}
x=a\cos t,\\
y=b\sin t,
\end{cases}
\qquad t\in[0,2\pi].
\]
一般地，设曲线的参数方程为
\[
\begin{cases}
x=x(t),\\
y=y(t),
\end{cases}
\qquad t\in[a,b].
\]
若 $x(t)$ 有反函数 $t=t(x)$，则可得复合函数
\[
y=y(t(x)).
\]
进一步设 $x(t)$ 和 $y(t)$ 在 $[a,b]$ 连续、可导，且 $x'(t)\neq0$．由复合函数求导法和反函数求导法得
\[
\frac{\mathrm{d}y}{\mathrm{d}x}=\frac{\mathrm{d}y}{\mathrm{d}t}\cdot\frac{\mathrm{d}t}{\mathrm{d}x}=\frac{y'(t)}{x'(t)}.
\]

\begin{example}\label{ex:4-3-5}
已知椭圆参数方程为
\[
\begin{cases}
x=a\cos t,\\
y=b\sin t,
\end{cases}
\qquad t\in[0,2\pi],
\]
求 $\dfrac{\mathrm{d}y}{\mathrm{d}x}$．

解 $\dfrac{\mathrm{d}y}{\mathrm{d}x}=\dfrac{y'(t)}{x'(t)}=\dfrac{b\cos t}{-a\sin t}=-\dfrac{b}{a}\cot t$．
\end{example}

\begin{example}\label{ex:4-3-6}
一轮子沿一直线滚动，轮子上一定点的轨迹曲线（见图~\ref{fig:4-8}）称为旋轮线，其参数方程为
\[
\begin{cases}
x=a(t-\sin t),\\
y=a(1-\cos t),
\end{cases}
\qquad 0\leqslant t\leqslant 2\pi.
\]
求出曲线上斜率为 $1$ 的切线．

\begin{figure}[htbp]
\centering
\includegraphics[width=0.55\textwidth]{fig4-8.png}
\figcap
\label{fig:4-8}
\end{figure}

解 旋轮线上任一点切线的斜率为
\[
y'_x=\frac{y'_t}{x'_t}=\frac{a\sin t}{a(1-\cos t)}=\cot\frac{t}{2}.
\]
令 $y'_x=1$，解得 $t=\dfrac{\pi}{2}$，它对应旋轮线上的点 $\left(a\left(\dfrac{\pi}{2}-1\right),a\right)$，故斜率为 $1$ 的切线为
\[
y-a=x-a\left(\frac{\pi}{2}-1\right),
\]
化简得
\[
x-y+a\left(2-\frac{\pi}{2}\right)=0.
\]
\end{example}

\begin{xiti}
\item 求下列隐函数的导数 $\dfrac{\mathrm{d}y}{\mathrm{d}x}$：
\begin{ziti}
\item $\dfrac{x^{2}}{a^{2}}+\dfrac{y^{2}}{b^{2}}=1$（$a$，$b$ 为常数）；
\item $y^{2}=2px$（$p$ 为常数）；
\item $x^{2}+xy+y^{2}=a^{2}$（$a$ 为常数）；
\item $x^{3}+y^{3}-xy=0$；
\item $y=x+\dfrac{1}{2}\sin y$；
\item $x^{\frac{2}{3}}+y^{\frac{2}{3}}=a^{\frac{2}{3}}$（$a$ 为常数）；
\item $y=\cos(x+y)$；
\item $y=x+\arctan y$；
\item $y=1-\ln(x+y)+e^{y}$；
\item $\arctan\dfrac{y}{x}=\ln\sqrt{x^{2}+y^{2}}$．
\end{ziti}

\item 求下列参数方程的导数：
\begin{ziti}
\item $\begin{cases} x=\dfrac{t}{1+t},\\ y=\dfrac{1-t}{1+t};\end{cases}$
\item $\begin{cases} x=\sin^{2}t,\\ y=\cos^{2}t;\end{cases}$
\item $\begin{cases} x=e^{2t}\cos^{2}t,\\ y=e^{2t}\sin^{2}t;\end{cases}$
\item $\begin{cases} x=a\left(\ln\tan\dfrac{t}{2}+\cos t\right),\\ y=a\sin t.\end{cases}$
\end{ziti}

\item 求函数 $y=y(x)$ 在指定点的导数：
\begin{ziti}
\item $y=\cos x+\dfrac{1}{2}\sin y$，$\left(\dfrac{\pi}{2},0\right)$；
\item $ye^{x}+\ln y=1$，$(0,1)$；
\item $\begin{cases} x=t-\sin t,\\ y=1-\cos t,\end{cases}$ 在 $t=\dfrac{\pi}{2},\pi$ 处；
\item $\begin{cases} x=1-t^{2},\\ y=t-t^{3},\end{cases}$ 在 $t=\dfrac{\sqrt{2}}{2}$，$\dfrac{\sqrt{3}}{3}$ 处．
\end{ziti}

\item 一圆锥形容器，深 $10\ \mathrm{m}$，上顶圆半径为 $4\ \mathrm{m}$．
\begin{ziti}
\item 灌入水时，求水的体积 $V$ 对水面高度 $h$ 的变化率；
\item 求体积 $V$ 对容器截面圆半径 $R$ 的变化率．
\end{ziti}

\item 设 $x=a\cos^{3}t$，$y=a\sin^{3}t$．
\begin{ziti}
\item 求 $y'(x)$；
\item 证明曲线的切线被坐标轴所截长度为一常数．
\end{ziti}

\item 证明：曲线 $\begin{cases} x=a(\cos t+t\sin t),\\ y=a(\sin t-t\cos t)\end{cases}$ 上任一点的法线到原点的距离恒等于 $a$．
\end{xiti}
